\section{Optimality and separation of the support}\label{sec:support}

We now consider the range that is not excluded by \cref{lem:initial-comparison}. The first-variation condition and conditional moment relations below follow Chi, Sun, and Wang~\cite[Observation 4.1 and Corollary 4.2]{CSW2026}.
Throughout this section, $k\ge6$, $x=\beta_k$, and $P$ is a maximizing distribution with $0<w<W_k$.
Thus $W_4<w<W_k$.
We retain the notation $a=kx$, $\delta=1-a=cx$, and $c=c_k$ from \cref{sec:comparison}.
For independent variables $X_i$ with distribution $P$, define
\begin{align}
 S_j&\coloneqq\sum_{i\in[j]}X_i,
 &F&\coloneqq F_k(P),
 &H_0&\coloneqq F_{k-1}(P),
 &p&\coloneqq1-F,
 &p_0&\coloneqq(1-x)^k=1-a^k.
 \label{eq:optimal-parameters}
\end{align}
Since the standard distributions attain $a^k$, we have $0<p\le p_0$; positivity follows also from $\E S_k=a<1$.
The estimates in \cref{lem:crossing-bounds} give
\begin{align}
 \frac{11}{24}+\frac{2}{25k}<c<\frac12,
 \qquad \frac{12}{13}<a<1,
 \qquad \frac13<p_0<e^{-1}.
 \label{eq:critical-numerical-bounds}
\end{align}

The next step uses the fact that $P$ maximizes $F_k$, rather than merely satisfying the mean constraint.
For a fixed value $z$ of the last variable, the probability of success is
\begin{align*}
 h(z)\coloneqq\Pr(S_{k-1}\ge1-z).
\end{align*}
We compare this function with an affine function whose average under $P$ is also $F$.

For $z\in[0,1]$, consider the signed perturbation
\begin{align*}
 P_\varepsilon\coloneqq P+\varepsilon\left[\delta_z-\frac zxP+\left(\frac zx-1\right)\delta_0\right].
\end{align*}
The expression in brackets has total mass and first moment zero.
Moreover, \cref{lem:normalization} gives $P(\{0\})\ge\delta>0$, so $P_\varepsilon$ is a probability distribution for all sufficiently small positive $\varepsilon$.
Differentiating the polynomial $F_k(P_\varepsilon)$ at $\varepsilon=0$ and using maximality gives
\begin{align*}
 0\ge \left.\frac{d}{d\varepsilon}F_k(P_\varepsilon)\right|_{\varepsilon=0}
 =k\left[h(z)-\frac zxF+\left(\frac zx-1\right)H_0\right].
\end{align*}
Consequently,
\begin{align}
 h(z)\le H_0+\frac{F-H_0}{x}z
 \qquad(z\in[0,1]).
 \label{eq:supporting-line}
\end{align}
The two sides have the same integral against $P$, because $\E h(X_1)=F$ and $\E X_1=x$.
Hence equality holds in \cref{eq:supporting-line} for $P$-almost every $z$.
We call such a value a \emph{contact point}.

At $z=1$, \cref{eq:supporting-line} gives $F-H_0\ge x(1-H_0)>0$, where $H_0<1$ follows from $(k-1)x<1$ and Markov's inequality.
Define
\begin{align}
 M&\coloneqq\frac{x(1-H_0)}{F-H_0},
 &D&\coloneqq M-x=\frac{xp}{F-H_0},
 &r&\coloneqq\E[1-S_k\mid S_k<1],
 &B&\coloneqq r-\delta.
 \label{eq:failure-parameters}
\end{align}
The preceding inequality and $p>0$ imply $x<M\le1$.
The affine function in \cref{eq:supporting-line} equals $1$ at $M$.
Since $h\le1$ and equality holds at almost every sampled value, $P([0,M])=1$.
At a contact point $z$, we therefore have
\begin{align}
 \Pr(S_{k-1}<1-z)=\frac{p(M-z)}D.
 \label{eq:contact-failure}
\end{align}
By conditioning on $X_1$, this identity determines the distribution of one summand when the sum fails:
\begin{align}
 \E[\varphi(X_1)\mid S_k<1]
 =\frac1D\E[(M-X_1)\varphi(X_1)]
 \label{eq:failure-tilt}
\end{align}
for every bounded Borel function $\varphi$.
In particular, failure gives a linear reweighting of the original distribution.

We record three consequences needed below.
Taking $\varphi(z)=z$ in \cref{eq:failure-tilt} gives
\begin{align*}
 1-r=\E[S_k\mid S_k<1]
 =\frac{k(Mx-\E X_1^2)}D.
\end{align*}
Since $M-x=D$ and $kx=1-\delta$, rearrangement yields
\begin{align}
 \Var(S_k)=k(\E X_1^2-x^2)=DB>0.
 \label{eq:unconditional-variance}
\end{align}
The strict inequality holds because $P$ has positive mean and positive mass at zero.
Thus $B>0$.
The identity $(1-y)_+-(y-1)_+=1-y$, together with \cref{eq:failure-parameters}, also gives
\begin{align}
 pr=\delta+\E(S_k-1)_+,
 \qquad 1-H_0=\frac{pM}D.
 \label{eq:failure-identities}
\end{align}

The first moment relation above is not sufficient to determine $P$.
We next compare the unconditional variance with the variance conditional on failure.
The following identity is Chi, Sun, and Wang~\cite[Lemma 4.3]{CSW2026}. It retains an additional nonnegative term; later, that term will also measure the separation between the two parts of the distribution.

\begin{lemma}\label{lem:covariance}
Let $v\coloneqq\Var(S_k\mid S_k<1)$, and let $X,X'$ be independent with distribution $P$.
Then
\begin{align}
 B(M-r)
 =v+\frac{k}{2p}\E\left[(X-X')\int_{X'}^X\bigl(h(y)-h(0)\bigr)\,dy\right].
 \label{eq:covariance-identity}
\end{align}
In particular, $r\le M$ and $v\le B(M-r)$.
\end{lemma}

\begin{proof}
Define $K(z)\coloneqq\E(1-S_{k-1}-z)_+$.
Independence and symmetry give
\begin{align*}
 \Cov(X,K(X))
 &=\E[(X_1-x)(1-S_k)_+]\\
 &=\frac1k\E[(S_k-kx)(1-S_k)_+].
\end{align*}
Conditional on $S_k<1$, the variable $1-S_k$ has mean $r$ and variance $v$.
As $S_k-kx=\delta-(1-S_k)$, the last expression equals
\begin{align}
 \Cov(X,K(X))=-\frac pk\bigl(v+rB\bigr).
 \label{eq:covariance-first}
\end{align}

On the other hand, symmetrizing the covariance gives
\begin{align*}
 \Cov(X,K(X))
 =\frac12\E[(X-X')(K(X)-K(X'))].
\end{align*}
The function $K$ is absolutely continuous and satisfies
\begin{align*}
 K(z)-K(z')=-\int_{z'}^z(1-h(y))\,dy.
\end{align*}
Separating the constant $1-h(0)=1-H_0$ inside this integral, and then using \cref{eq:unconditional-variance,eq:failure-identities}, gives
\begin{align*}
 \Cov(X,K(X))
 &=-(1-H_0)\Var(X)+\frac12\E\left[(X-X')\int_{X'}^X(h(y)-h(0))\,dy\right]\\
 &=-\frac{pMB}{k}+\frac12\E\left[(X-X')\int_{X'}^X(h(y)-h(0))\,dy\right].
\end{align*}
Comparing with \cref{eq:covariance-first} proves \cref{eq:covariance-identity}.
The function $h$ is nondecreasing on $[0,M]$, so the integrand after multiplication by $X-X'$ is nonnegative.
Since $B>0$, both conclusions follow.
\end{proof}

We use the interpolants from \cref{sec:comparison} to show that the difference $M-r$ in \cref{eq:covariance-identity} is small.
Choose the integer $t$ with $4\le t<k$ and $W_t\le w<W_{t+1}$.
Define
\begin{align}
 E^*&\coloneqq\frac{a^k}{k}\left(w+\sum_{j<k}b_je_j^*(w)\right),
 &V^*&\coloneqq\sum_{j<k}v_jf_j^*(w),
 &L&\coloneqq\frac{t(t-1)}{2k},\notag\\
 r_b&\coloneqq\frac{x(t+c)H(W_t)}{p_0},
 &r^*&\coloneqq\frac{\delta+E^*}{p_0},
 &M^*&\coloneqq x+\frac{xp_0}{a^k-V^*}.
 \label{eq:comparison-definitions}
\end{align}
Here, as before, sums indexed by $j<k$ run over $1\le j<k$.
The quantities $E^*$ and $V^*$ are the interpolated values of the excess and the $(k-1)$-variable tail.
The next lemma collects the estimates needed here from Chi, Sun, and Wang~\cite[Lemmas 4.6--4.9]{CSW2026}.
Its proof is deferred to \cref{sec:estimates}, where only scalar functions of $k,t,c,w$ are considered.

\begin{lemma}\label{lem:comparison-estimates}
With the quantities in \cref{eq:comparison-definitions}, we have
\begin{align}
 a^k>V^*,\qquad
 r_b>x,\qquad
 \frac{Lx}{(r_b-x)^2}<\frac9{10},\qquad
 r^*-\frac{19}{20}M^*>\frac x5.
 \label{eq:comparison-bounds}
\end{align}
Moreover, the multipliers in \cref{lem:interpolation} are nonnegative when its weights and costs are
$\kappa_j=Lv_j$ and $c_j=a^kb_j/k$.
\end{lemma}

The weights $v_j$ increase with $j$, since
\begin{align*}
 \frac{v_{j+1}}{v_j}=\frac{k(k-1-j)}{c(j+1)}>1
 \qquad(1\le j<k-1).
\end{align*}
Applying \cref{lem:interpolation} first with weights $v_j$ and zero costs, and then with weights $Lv_j$ and costs $a^kb_j/k$, gives
\begin{align}
 H_0\le V^*,
 \qquad
 \E(S_k-1)_+-LH_0\ge E^*-LV^*.
 \label{eq:actual-interpolated-comparison}
\end{align}
Indeed, the first assertion follows from \cref{eq:thinning-short-tail}.
For the second, subtract the weighted inequality from the common term $a^kw/k$ and use \cref{eq:thinning-excess}.

The next bound is Chi, Sun, and Wang~\cite[Lemma 4.4]{CSW2026}, expressed in the present notation.
\begin{lemma}\label{lem:mean-deficit}
The support endpoint $M$ and the conditional deficit $r$ satisfy
\begin{align}
 0\le M-r<\frac{D-3x}{20},
 \qquad D>3x.
 \label{eq:mean-deficit-proximity}
\end{align}
\end{lemma}

\begin{proof}
Use the representation $X_i=I_iY_i$, where the $I_i$ are independent Bernoulli variables with mean $a$ and the $Y_i$ are independent with distribution $Q$.
Conditional on the $Y_i$, if $R_k<1$, then $S_k<1$ and the expected deficit is $1-aR_k$.
Discarding the nonnegative contribution from $R_k\ge1$, we obtain
\begin{align*}
 pr=\E(1-S_k)_+
 &\ge\E[(1-aR_k)\1_{\{R_k<1\}}]\\
 &=\delta\Pr(R_k<1)+a\E(1-R_k)_+\\
 &\ge\delta H(w)+xw.
\end{align*}
For the second equality, we used $\E R_k=1$, which also implies $k\E(1-R_k)_+=w$.
The function $H$ is concave between consecutive nodes, and its values at the nodes increase.
It follows that $H(w)\ge H(W_t)$.
Since $W_t=tH(W_t)$ and $\delta=cx$, we have
\begin{align*}
 pr\ge x(t+c)H(W_t).
\end{align*}
Dividing by $p\le p_0$ gives $r\ge r_b$.

The definition of $M$, the inequalities $F\ge a^k$ and $H_0\le V^*$, and \cref{eq:comparison-bounds} give
\begin{align*}
 M=x+\frac{xp}{F-H_0}\le x+\frac{xp_0}{a^k-V^*}=M^*.
\end{align*}
Next, \cref{eq:failure-identities,eq:actual-interpolated-comparison} imply
\begin{align*}
 r+\frac{LM}{M-x}
 &=\frac{\delta+\E(S_k-1)_++L(1-H_0)}p\\
 &\ge\frac{\delta+E^*+L(1-V^*)}{p}\\
 &\ge\frac{\delta+E^*+L(1-V^*)}{p_0}
 =r^*+\frac{LM^*}{M^*-x}.
\end{align*}
The numerator in the third line is positive.
For the last equality, note that $M^*/(M^*-x)=(1-V^*)/p_0$.
Thus
\begin{align*}
 r^*-r
 &\le L\left(\frac{M}{M-x}-\frac{M^*}{M^*-x}\right)\\
 &=\frac{Lx(M^*-M)}{(M-x)(M^*-x)}
 \le\frac9{10}(M^*-M),
\end{align*}
where we used $M^*\ge M\ge r\ge r_b$ and \cref{eq:comparison-bounds}.
Consequently,
\begin{align*}
 r-\frac{19}{20}M
 &\ge r^*-\frac{19}{20}M^*+\frac1{20}(M^*-M)
 >\frac x5.
\end{align*}
Rearranging gives $M-r<(M-4x)/20=(D-3x)/20$.
Since \cref{lem:covariance} gives $M-r\ge0$, we also obtain $D>3x$.
\end{proof}

The next inequality applies to sums of independent centered variables with bounded ranges. This cubic argument replaces the normal-approximation steps in Chi, Sun, and Wang~\cite[Sections 5--6]{CSW2026}.
It will first show that the upper part of the support is nonempty, and then control the sum of the lower observations in \cref{sec:exclusion}.

\begin{lemma}\label{lem:cubic}
Let $Z_i$ be independent centered random variables, each supported on an interval of length at most $b>0$.
Write $Z\coloneqq\sum_iZ_i$ and $V\coloneqq\Var Z>0$.
For $\tau>0$ and $z\in\R$, set $Q\coloneqq\E[(Z+z)^2\1_{\{Z<\tau\}}]$.
Then
\begin{align}
 \left(V+\frac{(z-b)_+^2}{12}\right)Q
 \ge\frac{(\tau+z-2b)_+^3}{3}\E[Z\1_{\{Z\ge\tau\}}].
 \label{eq:cubic-test}
\end{align}
\end{lemma}

\begin{proof}
We first establish an inequality for a suitable test function.
Let $f$ be absolutely continuous, and let $G$ be nonnegative.
Suppose that $f'(y)\le G(w)$ whenever $|y-w|\le b$, with derivatives chosen arbitrarily at the finitely many corners of the function used below.
For each $i$, write $T_i\coloneqq Z-Z_i$, and take an independent copy $Z_i'$ of $Z_i$.
Since $\E Z_i=0$ and $Z_i$ is independent of $T_i$,
\begin{align*}
 \E[Z_if(Z)]
 &=\E[Z_i(f(T_i+Z_i)-f(T_i))]\\
 &=\E\left[Z_i^2\int_0^1 f'(T_i+\alpha Z_i)\,d\alpha\right].
\end{align*}
The support interval of $Z_i$ contains zero, because $Z_i$ is centered.
Thus $0$, $\alpha Z_i$, and $Z_i'$ lie in a common interval of length at most $b$.
The derivative assumption and independence give
\begin{align*}
 \E[Z_if(Z)]\le\E[Z_i^2G(T_i+Z_i')]
 =\E Z_i^2\,\E G(Z).
\end{align*}
Summing over $i$, we obtain
\begin{align}
 \E[Zf(Z)]\le V\E G(Z).
 \label{eq:test-function-bound}
\end{align}

If $L\coloneqq\tau+z-2b\le0$, then \cref{eq:cubic-test} is immediate.
Otherwise choose
\begin{align*}
 f(w)\coloneqq\frac13\min\{(w+z-b)_+^3,L^3\},
 \qquad
 G(w)\coloneqq(w+z)^2\1_{\{w<\tau\}}.
\end{align*}
Set $f'=0$ at the two corners.
If $f'(y)>0$, then $0<y+z-b<L$, so $y<\tau-b$.
For $|y-w|\le b$, it follows that $w<\tau$ and $w+z\ge y+z-b>0$.
Hence $f'(y)=(y+z-b)^2\le G(w)$, as required.

On $Z\ge\tau$, we have $f(Z)=L^3/3$, whereas on $0\le Z<\tau$, the product $Zf(Z)$ is nonnegative.
To control negative values, put $W\coloneqq Z+z$.
If $Z<0$ and $f(Z)>0$, then $b<W<z$, and
\begin{align*}
 -Zf(Z)
 &\le\frac{(z-W)(W-b)^3}{3}\\
 &\le\frac{(z-b)_+^2}{12}W^2.
\end{align*}
The last inequality uses $(z-W)(W-b)\le(z-b)^2/4$ and $(W-b)^2\le W^2$.
Therefore
\begin{align*}
 \E[Zf(Z)]
 \ge\frac{L^3}{3}\E[Z\1_{\{Z\ge\tau\}}]
 -\frac{(z-b)_+^2}{12}Q.
\end{align*}
Combining this with \cref{eq:test-function-bound}, where $\E G(Z)=Q$, proves the lemma.
\end{proof}

We can now separate the possible values of a maximizing variable.
The statement follows the support-separation argument of Chi, Sun, and Wang~\cite[Proposition 5.1]{CSW2026}, using \cref{lem:cubic} to show that the upper part is nonempty. The constants leave a strict inequality in both applications of the conditional variance bound.

\begin{lemma}\label{lem:support-gap}
The distribution $P$ gives full mass to
\begin{align*}
 [0,x+D/16]\ \cup\ (x+17D/20,M],
\end{align*}
and the second interval has positive mass.
\end{lemma}

\begin{proof}
For this proof, write $u\coloneqq x/D$ and $\mu\coloneqq(M-r)/D$.
By \cref{lem:mean-deficit}, $0<u<1/3$ and $0\le\mu<(1-3u)/20$.
Moreover,
\begin{align*}
 \frac BD=1+(1-c)u-\mu\le1+\frac{13u}{24}-\mu.
\end{align*}
The function $\mu(1+13u/24-\mu)$ increases for $0\le\mu\le(1-3u)/20$.
Applying \cref{lem:covariance} therefore gives
\begin{align}
 \frac v{D^2}
 \le\mu\frac BD
 \le\frac{(114+83u)(1-3u)}{2400}
 <\frac{19}{400}.
 \label{eq:small-conditional-variance}
\end{align}
For the last inequality, expand the numerator as $114-259u-249u^2$ and use $u>0$.

Let $W$ have the conditional distribution of $1-S_k$ given $S_k<1$.
Thus $\E W=r$ and $\Var W=v$.
If a contact point has the form $y=x+\xi D$ with $1/16\le\xi\le17/20$, then nonnegativity of the summands and \cref{eq:contact-failure} imply
\begin{align*}
 \Pr(W>y)
 =\frac{\Pr(S_k<1-y)}p
 \le\frac{\Pr(S_{k-1}<1-y)}p
 =1-\xi.
\end{align*}
Consequently $\Pr(W\le y)\ge\xi$.
Since $r-y=D(1-\mu-\xi)>0$, Cantelli's inequality gives
\begin{align*}
 \frac v{D^2}
 \ge\frac{\xi(1-\mu-\xi)^2}{1-\xi}
 \ge\frac{\xi(19/20-\xi)^2}{1-\xi}.
\end{align*}
On $[1/16,17/20]$, the derivative of the last function has the sign of $19/20-3\xi+2\xi^2$.
It first increases and then decreases, so its minimum occurs at an endpoint.
The two endpoint values are $5041/96000$ and $17/300$, both greater than $19/400$.
This contradicts \cref{eq:small-conditional-variance}.
Almost every value sampled from $P$ is a contact point, so the stated union has full mass.

Suppose now that the second interval has zero mass, and put $b_0\coloneqq x+D/16$.
Then every $X_i-x$ is supported on an interval of length at most $b_0$.
Apply \cref{lem:cubic} to
\begin{align*}
 Z=S_k-kx,
 \qquad \tau=\delta,
 \qquad z=B,
 \qquad b=b_0.
\end{align*}
The identities established above give
\begin{align*}
 Q=pv,\qquad V=DB,\qquad \E[Z\1_{\{Z\ge\delta\}}]=pB.
\end{align*}
For the first equality, note that $Z+B=S_k+r-1$ has conditional mean zero on failure.
The last equality follows by centering $Z$, whose conditional mean on failure is $-B$.
Furthermore,
\begin{align*}
 \frac{r-2b_0}{D}
 &=\frac78-u-\mu
 >\frac{33}{40}-\frac{17u}{20}
 >\frac{13}{24},\\
 \frac{(B-b_0)_+^2}{DB}
 &\le\frac BD<\frac43.
\end{align*}
After dividing \cref{eq:cubic-test} by $pDB$, these estimates yield
\begin{align*}
 \frac v{D^2}
 >\frac3{10}\left(\frac{13}{24}\right)^3
 =\frac{19}{400}+\frac{41}{230400}.
\end{align*}
This again contradicts \cref{eq:small-conditional-variance}, proving that the upper interval has positive mass.
\end{proof}
